\begin{table}[t]
\centering\scriptsize\begin{tabular}{lrr}
\toprule
r7 minus r6 & Difference & 95\% log interval \\
\midrule
Hit (pp) & +10.5363 & [+6.4624, +14.6212] \\
Early (pp) & -1.5225 & [-5.7790, +2.7510] \\
Missing (pp) & -4.3793 & [-9.3948, -0.2924] \\
Free (m) & +0.0420 & [-0.0005, +0.0697] \\
Distance (m) & -0.0794 & [-0.1874, -0.0093] \\
Recall (pp) & +8.4288 & [-0.7072, +23.9799] \\
\bottomrule
\end{tabular}

\caption{Completed joint r7 minus LiDAR r6 on five paired development logs. Geometry/hit gains coexist with unresolved free-space preservation.}
\label{tab:r7}
\end{table}
The final joint r7 run completes 30 epochs and 11,130 updates. Against matched LiDAR r6, it improves correct first hits by 10.536\,pp (95\% log interval $[6.462,14.621]$), with improvement in all five logs. Missing decreases 4.379\,pp ($[-9.395,-0.292]$), and one-way distance decreases 0.07940\,m ($[-0.18738,-0.00933]$). The joint pathway therefore has supported development gains in observed geometry and literal hits under this representation.

Physical preservation is not established. Early changes $-1.523$\,pp with interval $[-5.779,2.751]$; mean free intrusion increases 0.04202\,m with interval $[-0.00049,0.06974]$, worsening in four logs. Recall changes $+8.429$\,pp with interval $[-0.707,23.980]$. These intervals do not prove non-inferiority. The final r7 values are 49.604\% hit, 21.459\% early, 16.179\% missing, 0.27983\,m free intrusion, 0.09392\,m distance and 80.872\% recall. This result establishes a development increment only; the subsequent fixed external confirmation in Sec.~\ref{sec:external} fails to reproduce it. Free-space and surface-domain claims remain unresolved. The comparison changes native seeds, DPT features and build-depth auxiliary supervision together, with the same chart, physical objective and update budget. It does not attribute the gain to a single head or feature level.

A subsequent fixed-surface r7 audit uses the same 75 actors and 11,886 owned returns, with zero inference or optimization. Its ladder is 49.60\% $\to$ 61.97\% $\to$ 80.87\%; 12.36 percentage points still have an incorrect front surface with a correct rear intersection. Early triangles remain well-shaped. Build-supported components contribute 19.387 of the 21.459 early points, or 90.3\%. Removing unsupported components raises missing from 16.18\% to 25.11\%. This latest evidence retains the domain/first-surface problem despite r7's genuine development improvement.

r7 requires 20,249.42\,s (5.62 hours), including initial/final evaluation, with 10.224\,GiB peak CUDA allocation and 34.127\,GiB process RSS. It trains the 32.65M-parameter DPT and 1.71M-parameter query decoder. r6 takes 2,139.06\,s with 0.210\,GiB peak CUDA allocation. The comparison matches update count, not compute cost; these concurrent-system wall times are not exclusive GPU throughput.
